% =====================================================================
%  Dark Matter — front matter
%  Half-title, title page, copyright, dedication, epigraph, contents,
%  figures, tables, preface, how to use.
% =====================================================================

% ---- Half-title ---------------------------------------------------------
\pagestyle{plain}
\clearpage
\thispagestyle{empty}
\vspace*{1.15in}
{\sffamily\fontsize{22}{26}\selectfont\color{txt}\bookTitle\par}
\vspace{10pt}
{\sffamily\fontsize{11}{14}\selectfont\color{txtdim}\bookSubtitle\par}
\vfill
{\sffamily\small\color{txtdim}\bookAuthor\par}
\vspace{6pt}
{\sffamily\small\color{txtdim}\bookImprint\par}
\clearpage

% ---- Title page ---------------------------------------------------------
\clearpage
\thispagestyle{empty}
\vspace*{0.75in}
{\sffamily\fontsize{26}{30}\selectfont\bfseries\color{txt}\bookTitle\par}
\vspace{12pt}
{\sffamily\fontsize{12}{16}\selectfont\color{accent}\bookSubtitle\par}
\vspace{22pt}
\noindent{\color{accent}\rule{\textwidth}{1.8pt}}\par
\vspace{18pt}
{\sffamily\large\color{txt}\bookAuthor\par}
\vspace{8pt}
{\sffamily\color{txtdim}\bookImprint, publisher\par}
\vspace{4pt}
{\sffamily\color{txtdim}\bookEdition\par}
\vspace{18pt}
{\sffamily\small\color{txtdim}
Facts current as of \bookFactsDate.\par}
\vfill
\clearpage

% ---- Copyright page -----------------------------------------------------
\clearpage
\thispagestyle{plain}
\noindent{\sffamily\bfseries\color{txt}\bookTitle}\par
\noindent{\sffamily\small\color{txtdim}\bookSubtitle}\par
\vspace{14pt}
\noindent\rule{\textwidth}{0.4pt}\par
\vspace{10pt}
\noindent Copyright \textcopyright\ \bookCopyright\ \bookAuthor. All rights reserved.\par
\vspace{6pt}
\noindent \bookImprint, publisher.\par
\vspace{10pt}
\noindent\textbf{Edition:} \bookEdition.\par
\vspace{6pt}
\noindent\textbf{Subject headings:} Cosmology --- Dark matter; Astrophysics ---
Cosmology; Particle physics --- Dark matter; Direct detection (Astronomy);
Axions (Physics); Black holes (Astronomy).\par
\vspace{10pt}
\noindent\textbf{Licence.} This work is published under the Apache License,
Version 2.0. The complete licence text is available at
\href{https://www.apache.org/licenses/LICENSE-2.0}{\code{apache.org/licenses/LICENSE-2.0}}.
The book may be copied, modified and redistributed, including for commercial
purposes, provided that this notice and the attribution required by the licence
are retained.\par
\vspace{10pt}
\noindent\textbf{DOI:} \href{https://doi.org/\bookDOI}{\texttt{\bookDOI}}\par
\vspace{10pt}
\noindent\textbf{ISBN:} 978-1-0000-0000-0 (placeholder; to be assigned on
publication).\par
\vspace{10pt}
\noindent\textbf{Source materials.} The sources of this book, the build scripts,
the figures and the research ledger behind it are published at
\texttt{github.com/lordxmen2k/Dark-Matter-True-Origins}.\par
\vspace{10pt}
\noindent\textbf{What this book is.} This book is an assessment of a research
field at a stated date, written from primary sources. It reports what the
evidence establishes, what it does not, which explanations remain viable, and
what should be built, funded and decided next. It does not assert that dark
matter has been found, and it does not assert that any particular candidate is
correct. Where a claim could not be tied to a primary record, it is named as
unverified or dropped; the source ledger accompanying this book records that
process, including the claims that were dropped.\par
\vspace{10pt}
\noindent\textbf{Disclaimers.} The names of experiments, collaborations,
institutions and software projects mentioned in this book are the names of the
people and things whose work is being discussed. Their inclusion is not a
statement that they endorse, support or have verified anything described here.
Numerical results are quoted at the precision their source publications
support, and are date-stamped where they are likely to change.\par
\vspace{10pt}
\noindent\textbf{Facts current as of} \bookFactsDate. Statements that are likely
to change --- the status of the 248 keV candidate, the state of the exclusion
frontier, the status of any anomaly under replication --- are date-stamped
where they appear. Nothing in this book should be read as a statement about
events after that date.\par
\vspace{8pt}
\noindent{\color{accent}\rule{\textwidth}{0.4pt}}\par
\vspace{8pt}
{\small\color{txtdim}\sffamily
\noindent A note on the sources. Every identifier in the Notes and References was
resolved against the arXiv API, the Crossref bibliographic query API or
INSPIRE-HEP, and the resolved title was compared against the paper it was
intended to cite. Where a resolution failed, the claim was removed rather than
re-sourced from memory.\par}

% ---- Dedication ---------------------------------------------------------
\clearpage
\thispagestyle{empty}
\vspace*{2.5in}
  \begin{center}
    {\large\itshape\color{txt}
      To the people who kept the detectors cold through the winter,\\
      \hspace*{1em}and to the ones who read the appendix first.\par}
  \end{center}
  \vfill
\clearpage

% ---- Epigraph ------------------------------------------------------------
\clearpage
\thispagestyle{empty}
\vspace*{2.0in}
  \begin{center}
    \begin{minipage}{0.88\textwidth}
      \raggedright
      {\sffamily\color{accent}\rule{\textwidth}{2pt}}\par
      \vspace{14pt}
      {\large\itshape\color{txt}
      ``We have never touched it. Every number we hold about dark matter was
      inferred from its gravity, and the coupling we would need before that
      gravity could be called energy is about ten orders of magnitude out of
      reach.''\par}
      \vspace{12pt}
      {\raggedleft\sffamily\small\color{txtdim}
      The author\par}
    \end{minipage}
  \end{center}
  \vfill
\clearpage

% ---- Contents -----------------------------------------------------------
\clearpage
\pdfbookmark[0]{Contents}{contents}
\tableofcontents

% ---- List of figures ----------------------------------------------------
\clearpage
\pdfbookmark[0]{List of Figures}{list-of-figures}
\listoffigures

% ---- List of tables -----------------------------------------------------
\clearpage
\pdfbookmark[0]{List of Tables}{list-of-tables}
\listoftables

% ---- Preface ------------------------------------------------------------
\chapter*{Preface}
\addcontentsline{toc}{chapter}{Preface}
\markboth{Preface}{Preface}

\thispagestyle{plain}

Dark matter is the largest thing in the universe that nobody has ever held.
It accounts for roughly a quarter of all the energy in existence, it holds
galaxies together, it shapes the sky, and it has never produced a single
detectable event in any laboratory on Earth. It is also, by the standards of
physics, the best-evidenced hypothetical in the history of the subject: five
independent lines of evidence converge on the same number, and that number is
known to a few parts in a thousand.

That combination --- overwhelming evidence for the existence of something, and
no evidence at all about what it is --- is unusual, and it is the reason the
question deserves a book written for the present rather than for the archive.

The situation changed in a specific way over the year before this book was
written, and the change is what gives the assessment its shape. In September
2026 the LZ collaboration reported a single event in an energy window its
standard analysis had excluded, at a significance of 3.4$\sigma$ local and 2.6$\sigma$
global. Six interpretation papers appeared within four days. The event is
real, the data are public, and the analysis was not blinded. By the
standards the field has set for itself since the 3.5~keV line and the
21~cm anomaly, it is a candidate, not a discovery --- and the interesting
question is not whether it will survive, but what the field should do while
waiting to find out.

This book is written on the assumption that the next decade will not settle
the question by one dramatic event. It will settle it, if it settles it at
all, by a portfolio of experiments, astrophysical surveys and cosmological
measurements whose results are individually interpretable and collectively
decisive. The book's purpose is to say what that portfolio should contain,
what each experiment would and would not prove, and which candidates are
already known to be wrong.

Four commitments govern the writing.

The first is \textbf{that the evidence is read, not recalled}. Every number in
this book comes from a primary publication that was retrieved and checked
while the book was written, and every identifier was resolved against a
bibliographic registry rather than typed from memory. Where that process
failed, the claim is marked or removed. The source ledger that accompanies the
book is the record of that process, including its corrections.

The second is \textbf{that null results are treated as results}. The exclusion
curves of the last five years are the strongest constraints in the field, and
the most common public error is to read them as evidence that dark matter
does not exist or that the search has failed. They are neither. They are a
map of where the answer is not.

The third is \textbf{that the question ``have we already tapped it?'' is taken
seriously in each of its senses, including the ones that turn out to have
interesting answers}. We have used dark matter as a gravitational probe for
ninety years. We have captured a private supply of it in the Solar System,
where the local density is enhanced by orders of magnitude. We have, on
paper, harvested it. The book's claim is that the interesting senses of the
question have positive answers and the particle sense does not, and that the
distance between those two facts is the distance the field has to close.

The fourth is \textbf{that the reader is told what to do with the book}. Part
VI is not a summary. It contains a decision framework, a candidate-by-candidate
criterion set with cost, timescale and kill conditions, a statement of what a
discovery would and would not prove, and a roadmap to 2040 that says what
should be running when.

Readers who want the state of play and the decision framework can read
Part~III and Part~VI. Readers who want the physics can read Part~I and
Part~II. Readers who want the origins question --- the part of the title that
is hardest to answer honestly --- should read Part~IV and Chapter~\ref{ch:harvest},
which explains why the origins question is downstream of the experimental
programme rather than upstream of it.

% ---- How to use this book -----------------------------------------------
\chapter*{How to Use This Book}
\addcontentsline{toc}{chapter}{How to Use This Book}
\markboth{How to Use This Book}{How to Use This Book}

\thispagestyle{plain}

\begin{takeaway}{The shape of the book}
\begin{itemize}
\item \textbf{Part~\ref{part:budget}} fixes what is known: the energy budget, the
five independent handles on the missing mass, the failure of baryons to supply
it, and the small-scale structure that standard dark matter struggles to
produce.
\item \textbf{Part~\ref{part:candidates}} takes the candidate classes one at a
time. For each: what the idea is, what it predicts, what it has cost the field
so far, what would confirm it, and what would exclude it.
\item \textbf{Part~\ref{part:frontier}} is the state of the field in 2026: the
exclusion frontier and the neutrino fog, the 248~keV candidate, the anomaly
record, the cosmological probes, and the tensions that cosmology itself has
opened.
\item \textbf{Part~\ref{part:origins}} asks where the dark matter came from.
The relic density is a constraint, not an explanation, and the chapter on
production mechanisms shows which mechanisms explain it and which merely
reproduce it.
\item \textbf{Part~\ref{part:tapped}} answers the title question in five
senses, with the arithmetic attached to each.
\item \textbf{Part~\ref{part:decision}} turns the assessment into a programme:
why the answer must be a portfolio, what each candidate's decision criteria
are, what a discovery would prove, and what should run when.
\end{itemize}
\end{takeaway}

Three readings are possible, and each is coherent.

\emph{If you want to know where the field stands.} Read Part~III. It is the
compressed version of the book: where the frontier is, what the one live
candidate is, and what the record of non-replications implies for how anyone
should behave in 2027.

\emph{If you want the physics.} Read Part~I, then Part~II. No background beyond
undergraduate physics and astronomy is assumed; every result the argument
depends on is stated, sourced and explained where the statement alone is not
enough.

\emph{If you have to decide something.} Read Chapter~\ref{ch:criteria} and
Chapter~\ref{ch:roadmap}. The first is a set of decision criteria with kill
conditions attached to each candidate; the second is a schedule with the
experiments that should occupy it.

\begin{mythbox}{Five things this book will not do}
\begin{itemize}
\item It will not announce a discovery. Nothing in the evidence assembled here
meets the standard the field has adopted for one, and Chapter~\ref{ch:event248}
says exactly what the 248~keV event is and is not.
\item It will not declare dark matter undetectable. The fog is a floor for one
class of search, not a verdict on the field; Chapter~\ref{ch:direct} sets out
which channels remain open and why.
\item It will not rank the candidates by popularity. Popularity in this field
has been a poor predictor of survival for twenty years.
\item It will not treat a failed replication as evidence of fraud. The anomaly
record in Chapter~\ref{ch:indirect} is a record of statistical and modelling
error, not of misconduct.
\item It will not discuss how the book was produced. The methods behind it are
in the accompanying source materials, not in the text.
\end{itemize}
\end{mythbox}

\vspace{8pt}
\noindent{\small\color{txtdim}\sffamily
Conventions. Mathematical statements are given in ordinary notation. Cross
sections are quoted in the units their source publications use, with the local
density assumption stated wherever a limit is compared with another. Where the
book says that something was verified, it means that the identifier was
resolved against a bibliographic registry and the resolved title matched the
paper intended --- not that it was recalled. Where the book says
``unverified'', it means exactly that, and the sentence in question says what
would be needed to verify it.\par}
